% =====================================================================
%  Seksjon: Den ytre frekvensstøtte-droopen i vindturbinmodellen.
%  Kilde:   src/tops_openfast/dyn_models/windturbine_tower.py
%             _droop_command(), P_ref_components(), set_grid_frequency_hz()
%  Driver:  casestudies/dyn_sim/test_WT_LEOGO_tower_sim.py  (feed_grid_frequency)
%  Preamble: amsmath, amssymb, graphicx, booktabs, siunitx
%           + \DeclareSIUnit{\pu}{pu}   (per-unit, used by \SI{}{\pu})
% =====================================================================

\section{The Outer Wind-Turbine Frequency-Support Droop}
\label{sec:outer-droop}

To place the tower-mode interaction in an operationally realistic setting, the
reduced turbine of \S\ref{sec:simplified-tower-model} carries an optional
\emph{frequency-support droop}: on an islanded platform grid such as LEOGO the
wind turbine is expected to help arrest frequency excursions after a disturbance,
the same service the gas-turbine governors provide. This section defines that
outer droop law; whether it is actually needed, given that the converter already
contains a droop, is the subject of \S\ref{sec:outer-droop-necessity}.

\subsection{Control law}
\label{subsec:od-law}

The droop is implemented in \texttt{WindTurbineTower.\_droop\_command}. The
turbine is first \emph{de-loaded}: instead of tracking the maximum-power-tracking
(MPT) point \(P_{\mathrm{avail}}\), it holds a reduced base reference
\begin{equation}
  P_{\mathrm{base}} = \operatorname{clip}\!\big(
    P_{\mathrm{avail}} - P_{\mathrm{hr}},\;0,\;P_{\mathrm{avail}}\big),
  \label{eq:od-base}
\end{equation}
keeping a headroom reserve \(P_{\mathrm{hr}}\) (parameter
\(\texttt{headroom\_pu}\), on the UIC base) available for upward support. On top
of the base it adds a proportional response to the measured grid-frequency
deviation,
\begin{equation}
  \Delta P = K_{\mathrm{droop}}\,\big(f_{\mathrm{nom}} - f_{\mathrm{grid}}\big),
  \qquad
  P_{\mathrm{ref}} = \operatorname{clip}\!\big(
    P_{\mathrm{base}} + \Delta P,\;0,\;P_{\mathrm{avail}}\big),
  \label{eq:od-delta}
\end{equation}
so that a frequency dip (\(f_{\mathrm{grid}} < f_{\mathrm{nom}}\)) raises the
power reference and injects extra active power into the grid, up to the headroom
limit. The gain \(K_{\mathrm{droop}}\) is \(\texttt{k\_droop\_pu\_per\_hz}\) and
\(f_{\mathrm{nom}} = \SI{50}{\hertz}\). The droop is toggled with
\(\texttt{droop\_enable}\); when it is off, \(\Delta P = 0\) and the turbine
simply holds the de-loaded base \eqref{eq:od-base}.

An important detail for the experiments is that the headroom \eqref{eq:od-base}
is applied \emph{whether or not} the droop is enabled. Both the droop-off and
droop-on runs therefore share the same de-loaded operating point, and their
difference isolates the droop \emph{response} \eqref{eq:od-delta} alone, with no
confounding change in the base loading.

\subsection{Measured grid frequency}
\label{subsec:od-freq}

The frequency the droop reacts to is supplied from the network side of the
co-simulation through \texttt{set\_grid\_frequency\_hz}. In the LEOGO driver the
platform is dominated by the three coherent gas-turbine synchronous machines on
the main bus, so the frequency is taken as the centre-of-inertia (mean
generator-speed) frequency
\begin{equation}
  f_{\mathrm{grid}} = f_n\big(1 + \overline{\Delta\omega}_{\mathrm{gen}}\big),
  \label{eq:od-fgrid}
\end{equation}
with \(f_n = \SI{50}{\hertz}\) and \(\overline{\Delta\omega}_{\mathrm{gen}}\) the
mean per-unit generator-speed deviation. Because the generator speed is a state,
\(f_{\mathrm{grid}}\) is available before the algebraic solve at each integration
stage and is fed to the turbine, which reads it when it forms \(P_{\mathrm{ref}}\)
(\texttt{feed\_grid\_frequency} in the driver). This closes an outer,
network-measured loop around the turbine, distinct from the converter's own local
droop (\S\ref{sec:outer-droop-necessity}).

\subsection{Parameters and defaults}
\label{subsec:od-params}

Table~\ref{tab:od-params} lists the droop parameters. In the headline scenario of
\S\ref{subsec:nt-tradeoff} the droop uses \(K_{\mathrm{droop}} =
\SI{0.75}{\pu\per\hertz}\) with a \SI{10}{\percent} headroom reserve; these are
the values for which the frequency-support benefit and the tower-loading penalty
are reported.

\begin{table}[t]
  \centering
  \caption{Outer frequency-support droop parameters
  (\texttt{WindTurbineTower}).}
  \label{tab:od-params}
  \begin{tabular}{lll}
    \toprule
    Parameter & Symbol & Value \\
    \midrule
    Nominal frequency        & \(f_{\mathrm{nom}}\)      & \SI{50}{\hertz} \\
    Droop gain               & \(K_{\mathrm{droop}}\)    & \SI{0.75}{\pu\per\hertz} \\
    Headroom reserve         & \(P_{\mathrm{hr}}\)       & \SI{0.10}{\pu} (UIC base) \\
    Enable flag              & \(\texttt{droop\_enable}\) & \(\{0,1\}\) \\
    \bottomrule
  \end{tabular}
\end{table}
